% Introduction: motivation, VVUQ principles, thesis, contributions.
% The Introduction precedes the Table of Contents.

Physical artificial intelligence is entering oncology clinical trials along
three converging fronts. The first is in silico simulation, where digital twins
and mechanistic models stand in for cohorts of patients so that dose, schedule,
and endpoint decisions can be explored before a single person is enrolled. The
second is autonomous code generation, where large language model agents now
write the analysis code, the simulation harness, and even the statistical
programming that a trial depends on. The third, and the most consequential, is
the gradual arrival of physical procedures under machine control, from
laboratory automation to the robotic surgery and combined medicine pilots that a
2030 oncology program may rely on. 

Across all three fronts the asymmetry of cost
is the same and it is severe: an error that is caught before a deliverable ships
is a line in a log file, while an error that ships into a simulation that informs
a protocol, or into a procedure that touches a patient, is a clinical and a
regulatory event. The faster and the more autonomous the generation becomes, the
more this asymmetry matters, because speed without a increase in
assurance simply ships mistakes sooner.

The discipline that addresses this asymmetry is verification, validation, and
uncertainty quantification, abbreviated VVUQ, and it is now well established for
digital twins in precision medicine~\cite{sel2025vvuq, moradi2026twins} and for
the physiologically based pharmacokinetic models that underpin oncology dosing
~\cite{zaid2025pbpk}. In the specific context of an oncology trial deliverable,
the three dimensions can be stated plainly. Verification asks ``Did we build the
code correctly?'', that is, does the software implement its specification, pass
its structural and schema checks, and behave as intended.

Validation asks ``Did we build the right model?'', that is, does the deliverable
agree with a trusted reference or ground truth to within an acceptable error, and
has a qualified human reviewed and signed off on it. Uncertainty quantification
asks ``How confident are we in the results?'', that is, when the work is repeated
under controlled perturbation, how dispersed are the outputs and is that
dispersion inside an agreed bound. These three questions are the substance of the
credibility frameworks that regulators and standards bodies have built for
computational evidence in medical contexts~\cite{asmevv40, fda2023credibility}.


The LLM VVUQ process
must be more substantial than the automated code generation it gates, and being
more substantial is precisely what makes autonomous oncology trial deliverables
faster, less expensive, and more rigorous than conventional verification. Modern agents can emit a complete, well formed deliverable
in well under a millisecond, and recent systems show that deterministic and
agentic code generation for clinical workflows is increasingly
reliable~\cite{trooskens2026compiled, carlisle2026pipeline, yan2026clinagent}.
The hard, decision bearing work is not producing that candidate; it is the gate
that decides whether the candidate may ship, the thresholds that gate enforces,
and the escalation policy that routes a doubtful case to a human. 

When that
assurance layer is automated, exercised across its full decision surface, and
held to a deliberately higher standard than the generator, the result is a loop
that is both quick and trustworthy. The platform analyzed here, the
\texttt{cancer-automated} repository, is the realization of that
loop~\cite{repo-cancer-automated}, and it continues a line of the author's
physical AI oncology trial automation that includes a fully automated sponsor
platform~\cite{kawchak2026sponsor} and a metastatic pancreatic cancer trial
established under an explicit code, VVUQ, and playbook
pattern~\cite{kawchak2025qsppdac}.

This paper makes four contributions, each developed in a later section and
referenced here so the argument reads as one connected piece. First, it presents
a single repeatable pipeline of five established methods, four producing stages
and one assurance gate, packaged so that one daily deliverable runs them in order
(\autoref{sec:methods}). Second, it presents a VVUQ gate that is held to a higher
standard than the generation it gates and is exercised across its full accept,
block, and escalate surface, which is the core empirical result of the work
(\autoref{sec:results}). Third, it gives a deliberately honest account of what
actually ran, what was approximated, and what could not be run in the execution
environment, because a thesis built on assurance is only credible if its own gaps
are reported plainly (\autoref{sec:limitations}). Fourth, it carries a Stage 2
deployment reference, a 2030 pancreatic cancer analog combining a robotic Whipple
procedure with an adjuvant medicine regimen (\autoref{sec:results} and
\autoref{sec:discussion}).

The surrounding literature motivates each half of this argument. On the assurance
side, the survey of VVUQ for precision medicine digital twins frames the
verification, validation, and uncertainty vocabulary used here~\cite{sel2025vvuq},
the image to physics to digital twin pathway shows where personalized cancer
models are heading~\cite{moradi2026twins}, and the application of VVUQ to
physiologically based pharmacokinetic models for theranostics is a direct oncology
analog of the dosing decisions a trial must defend~\cite{zaid2025pbpk}. 

The
credibility of such computational evidence is governed by established frameworks,
in particular the risk informed ASME V\&V 40 standard and the FDA guidance on
assessing credibility of modeling and simulation in device
submissions~\cite{asmevv40, fda2023credibility}, with good clinical practice and
medical device software life cycle expectations layered on
top~\cite{iche6r3, iec62304}. On the generation side, deterministic and agentic
code generation for clinical workflows is advancing quickly, from compiled and
deterministic generation~\cite{trooskens2026compiled} to validated generative
extraction pipelines~\cite{carlisle2026pipeline} and layered autonomous
statistical programming~\cite{yan2026clinagent}. The contribution here is not to
improve either half in isolation but to insist on the ordering between them, and
to do so at a moment when regulators are themselves moving toward faster, real
time clinical evidence~\cite{fda2026realtime}.

This work also stands on a substantial body of prior evidence. The repository's
input corpus references many distinct external DOIs spanning VVUQ theory, digital
twins, and in silico oncology trials, and the platform itself continues a
sequence of the author's physical AI oncology trial systems, from autonomous
glioblastoma meta analyses with explicit verification~\cite{kawchak2025glioblastoma}
and a large language model code generation benchmark on adverse event
data~\cite{kawchak2025faers}, to federated trial unification~\cite{kawchak2026federated}
and accelerated patient prediction~\cite{kawchak2026prediction}. 

Finally, this work positions the platform as the next step beyond the prior
humanoid instruction to code to execution to paper workflow developed in the
author's Clinical-AI-Demos project. That earlier workflow proved that an agent
can carry a task from a bracketed instruction through code, execution, and a
written paper without a human in the inner loop. What it left open, and what this
platform strengthens, is the assurance around that loop and the discipline of
per commit automation: a VVUQ gate held above generation, triple run
uncertainty, lossless autochunking under a strict file size cap, see \autoref{figtreemap}, and a record
where every section is its own commit pushed in real time. The sources that
ground each part of the argument are collected in \autoref{tab:intro_sources};
the prose that follows draws on them in the author's own words rather than
restating them.

\begin{table}[H]
\caption{Introduction source map (abbreviated paths under the repository root).}
\label{tab:intro_sources}
\centering
\begin{tabularx}{\textwidth}{L{6cm} >{\raggedright\arraybackslash}X}
\toprule
\textbf{Source (under repo root)} & \textbf{Use for} \\
\midrule
  \texttt{vvuq/} : README, 4 mods. & Three VVUQ dimensions, gate principle, stated plainly. \\
  \texttt{execution/README.md} : Thesis & Thesis wording and assurance vs.\ generation argument. \\
  \texttt{inputs/Paper-VVUQ-1, -2} : README and chunk 01 & Background on VVUQ for in silico trials and the broad DOI input corpus. \\
  \texttt{inputs/VVUQ-Research-1, -2} : README and chunk 1 & Additional VVUQ and digital twin framing for oncology. \\
  Clinical-AI-Demos 07-humanoid & The prior workflow this paper extends. \\
\bottomrule
\end{tabularx}
\end{table}

\vspace*{-1.25cm}

\begin{figure}[H]
    \hspace{-0.8cm}%
      {\includegraphics[width=1.1\linewidth, trim={0 25 0 0}, clip]{Images/06-test-coverage-treemap.png}}
    \vspace{-0.5cm}
    \caption{Test coverage treemap of the 51 passing tests across eight modules, each sized to test count.}
    \label{figtreemap}
\end{figure}